\section{Aspetti Implementativi}

Gli aspetti implementativi descrivono come l'architettura viene tradotta in file e responsabilita' concrete. Il progetto usa una struttura semplice, coerente con Next.js App Router: le pagine stanno in \texttt{src/app}, i componenti riutilizzabili in \texttt{src/components}, le funzioni di supporto in \texttt{src/lib} e la protezione della navigazione in \texttt{src/proxy.ts}.

\subsection{Struttura del progetto}

\begin{figure}[H]
\centering
\resizebox{0.98\textwidth}{!}{
\begin{tikzpicture}[
    root/.style={rectangle, draw=aliceblue, fill=aliceblue!12, thick, rounded corners=4pt, minimum width=3.4cm, minimum height=0.9cm, align=center, font=\small\bfseries},
    folder/.style={rectangle, draw=alicegreen, fill=alicegreen!10, thick, rounded corners=3pt, minimum width=3.5cm, minimum height=0.85cm, align=center, font=\scriptsize\bfseries},
    file/.style={rectangle, draw=gray, fill=gray!8, rounded corners=3pt, minimum width=3.5cm, minimum height=0.75cm, align=center, font=\scriptsize},
    line/.style={thick, color=aliceblue},
]
\node[root] (src) at (0, 0) {\texttt{src}};

\node[folder] (app) at (-5, -1.6) {\texttt{app}\\route e pagine};
\node[folder] (components) at (0, -1.6) {\texttt{components}\\UI riusabile};
\node[folder] (lib) at (5, -1.6) {\texttt{lib}\\helper};
\draw[line] (src) -- (app);
\draw[line] (src) -- (components);
\draw[line] (src) -- (lib);

\node[file] (op) at (-7.2, -3.2) {\texttt{operatore}\\dashboard e actions};
\node[file] (pz) at (-5, -3.2) {\texttt{paziente}\\home, umore, esercizi, farmaci};
\node[file] (api) at (-2.8, -3.2) {\texttt{api/ai}\\route Groq};
\draw[line] (app) -- (op);
\draw[line] (app) -- (pz);
\draw[line] (app) -- (api);

\node[file] (dash) at (-1.2, -3.2) {\texttt{DashboardContent}\\filtri e lista};
\node[file] (card) at (0.9, -3.2) {\texttt{PatientCard}\\riepilogo};
\node[file] (detail) at (3.0, -3.2) {\texttt{PazienteDettaglio}\\tab e gestione};
\draw[line] (components) -- (dash);
\draw[line] (components) -- (card);
\draw[line] (components) -- (detail);

\node[file] (date) at (4.0, -3.2) {\texttt{date.ts}};
\node[file] (form) at (5.3, -3.2) {\texttt{form.ts}};
\node[file] (metrics) at (6.6, -3.2) {\texttt{metriche.ts}};
\draw[line] (lib) -- (date);
\draw[line] (lib) -- (form);
\draw[line] (lib) -- (metrics);

\node[file] (proxy) at (0, -4.8) {\texttt{src/proxy.ts}\\sessione e redirect};
\draw[line] (src) -- (proxy);
\end{tikzpicture}
}
\caption{Struttura implementativa principale del progetto}
\label{fig:struttura-progetto}
\end{figure}

\subsection{Pagine e route}

Le route principali sono separate per ruolo. L'area operatore contiene login, registrazione, dashboard e dettaglio paziente. L'area paziente contiene login, registrazione, home e le tre funzioni giornaliere: umore, esercizi e farmaci. La route \texttt{/api/ai/suggerimento-operatore} non e' una pagina, ma un endpoint server-side richiamato dal dettaglio paziente.

\begin{table}[H]
\centering
\renewcommand{\arraystretch}{1.25}
\footnotesize
\begin{tabularx}{\textwidth}{|p{5.6cm}|X|}
\hline
\textbf{File} & \textbf{Responsabilita'} \\
\hline
\texttt{src/app/operatore/dashboard/page.tsx} & Carica l'operatore, i pazienti assegnati e le metriche di riepilogo. \\
\hline
\texttt{src/app/operatore/dashboard/pazienti/[id]/page.tsx} & Verifica che il paziente appartenga all'operatore e prepara dati di dettaglio. \\
\hline
\texttt{src/app/paziente/home/page.tsx} & Mostra la home semplificata del paziente. \\
\hline
\texttt{src/app/paziente/umore/page.tsx} & Salva o aggiorna l'umore giornaliero. \\
\hline
\texttt{src/app/paziente/esercizi/page.tsx} & Gestisce l'esecuzione guidata e il log degli esercizi. \\
\hline
\texttt{src/app/paziente/farmaci/page.tsx} & Mostra farmaci per fascia e registra l'assunzione. \\
\hline
\end{tabularx}
\caption{Pagine principali}
\label{tab:pagine-implementate}
\end{table}

\subsection{Server Actions e accesso ai dati}

Le Server Actions gestiscono le operazioni che modificano dati: registrazione, login, logout, aggiornamento del paziente, aggiunta o rimozione di farmaci ed esercizi. La scelta delle Server Actions evita controller REST manuali per ogni form e mantiene la logica lato server.

\begin{table}[H]
\centering
\renewcommand{\arraystretch}{1.25}
\begin{tabularx}{\textwidth}{|p{5.8cm}|X|}
\hline
\textbf{File} & \textbf{Operazioni principali} \\
\hline
\texttt{src/app/operatore/actions.ts} & Login, registrazione e logout operatore. \\
\hline
\texttt{src/app/paziente/actions.ts} & Login, registrazione e logout paziente. \\
\hline
\texttt{src/app/operatore/dashboard/pazienti/actions.ts} & Update paziente, create/delete farmaco, create/delete esercizio. \\
\hline
\end{tabularx}
\caption{Server Actions del progetto}
\label{tab:server-actions}
\end{table}

\subsection{Componenti client}

I componenti client sono usati dove serve interazione nel browser. \texttt{DashboardContent.tsx} gestisce ricerca e filtri nella dashboard; \texttt{PatientCard.tsx} visualizza un paziente in forma sintetica; \texttt{PazienteDettaglio.tsx} gestisce tab, modali, form, storico e richiesta del suggerimento AI.

\subsection{Helper di supporto}

Durante il refactoring alcune funzioni ripetute sono state spostate in \texttt{src/lib}. \texttt{date.ts} gestisce la data corrente, \texttt{form.ts} centralizza la lettura dei valori da \texttt{FormData}, \texttt{metriche.ts} calcola percentuali e \texttt{paziente-client.ts} recupera l'id del paziente associato all'utente autenticato.

\subsection{Suggerimento AI}

Il suggerimento AI e' isolato nella route \texttt{src/app/api/ai/suggerimento-operatore/route.ts}. Il client invia dati sintetici e non identificativi; la route usa la chiave \texttt{GROQ\_API\_KEY} lato server, chiama Groq tramite \texttt{fetch} e restituisce una sola frase per l'operatore. Questa separazione evita di inserire la chiave nel browser e mantiene la logica AI fuori dal componente React.

\subsection{Sicurezza}

La sicurezza si basa su Supabase Auth, proxy Next.js e Row Level Security. Il proxy protegge le pagine private e gestisce i redirect; le pagine server controllano che l'operatore acceda solo ai propri pazienti; le policy RLS nel database impediscono letture o modifiche non coerenti con il ruolo.

\newpage
